\paragraph{Making the Kicked-Back Phase Observable}\label{ex:making-the-kicked-back-phase-observable}

In the previous section we discovered the central result of phase kickback:
\begin{equation*}
    CU\ket{+}\ket{v} = \left(\frac{\ket{0}+e^{i\phi}\ket{1}}{\sqrt{2}}\right) \otimes \ket{v}
\end{equation*}
This is already a remarkable result: the eigenvalue phase has become a \textbf{relative phase} on the control qubit. However, a natural question immediately arises: \emph{\textbf{can we measure this phase directly?}} Surprisingly, the answer is \textbf{no}. This section explains why the kicked-back phase is \textbf{hidden from a direct measurement}, and why quantum algorithms always perform an additional interference step before measuring.

\highspace
\begin{flushleft}
    \textcolor{Green3}{\faIcon[regular]{bookmark} \textbf{The state after phase kickback}}
\end{flushleft}
Recall the final state after applying the controlled unitary:
\begin{equation*}
    \ket{\psi} = CU\ket{+}\ket{v} = \left(\frac{\ket{0}+e^{i\phi}\ket{1}}{\sqrt{2}}\right) \otimes \ket{v}
\end{equation*}
Since the target qubit is unchanged, we only need to examine the control qubit:
\begin{equation*}
    \ket{c} = \frac{\ket{0}+e^{i\phi}\ket{1}}{\sqrt{2}}
\end{equation*}
Everything we want to know is encoded in the phase $e^{i\phi}$. The question is therefore: \emph{\textbf{if we simply measure this qubit, can we recover $\phi$?}}

\highspace
\begin{flushleft}
    \textcolor{Red2}{\faIcon{times-circle} \textbf{Measuring in the computational basis}}
\end{flushleft}
Suppose we perform a standard measurement in the computational basis:
\begin{equation*}
    \left\{\ket{0}, \ket{1}\right\}
\end{equation*}
The amplitudes are:
\begin{equation*}
    \ket{\alpha} = \frac{1}{\sqrt{2}}, \quad \ket{\beta} = \frac{e^{i\phi}}{\sqrt{2}}
\end{equation*}
Measurement probabilities are obtained from the squared moduli:
\begin{itemize}
    \item Measuring $0$:
    \begin{equation*}
        P(0) = |\alpha|^2 = \left|\frac{1}{\sqrt{2}}\right|^2 = \frac{1}{2}
    \end{equation*}


    \item Measuring $1$ is a little more interesting:
    \begin{equation*}
        P(1) = |\beta|^2 = \left|\frac{e^{i\phi}}{\sqrt{2}}\right|^2
    \end{equation*}
    Using:
    \begin{equation*}
        \left|z\right|^2 = z^* z
    \end{equation*}
    Where $z^*$ is the complex conjugate of $z$, we have:
    \begin{align*}
        P(1) &= \left(\frac{e^{i\phi}}{\sqrt{2}}\right)^* \left(\frac{e^{i\phi}}{\sqrt{2}}\right) \\[.3em]
        &= \left(\frac{e^{-i\phi}}{\sqrt{2}}\right) \left(\frac{e^{i\phi}}{\sqrt{2}}\right) \\[.3em]
        &= \frac{e^{-i\phi} e^{i\phi}}{2} \\[.3em]
        &= \frac{e^{-i\phi + i\phi}}{2} \\[.3em]
        &= \frac{e^{0}}{2} \\[.3em]
        &= \frac{1}{2}
    \end{align*}
\end{itemize}
Therefore, the measurement probabilities are:
\begin{equation*}
    P(0) = P(1) = \frac{1}{2}
\end{equation*}
This means that the measurement outcome is completely independent of the phase $\phi$. So a \hl{computational-basis} measurement \textbf{cannot distinguish any of these states}. The reason lies in how quantum measurements work: \hl{a measurement does \textbf{not} use the complex amplitudes directly, instead it computes the squared modulus of the amplitudes.} Taking the modulus removes every complex phase factor, and therefore the measurement outcome is completely independent of the phase $\phi$.

\highspace
\begin{flushleft}
    \textcolor{DarkOrange3}{\faIcon{balance-scale} \textbf{Phase vs Amplitude Information}}
\end{flushleft}
A quantum state contains two different kinds of information:
\begin{enumerate}
    \item \textbf{Amplitude}: Determines the probability of measuring a particular state. For example:
    \begin{equation*}
        \frac{\sqrt{3}}{2}\ket{0} + \frac{1}{2}\ket{1}
    \end{equation*}
    Has:
    \begin{equation*}
        P(0) = \left|\frac{\sqrt{3}}{2}\right|^2 = \frac{3}{4}, \quad P(1) = \left|\frac{1}{2}\right|^2 = \frac{1}{4}
    \end{equation*}
    These probabilities are \emph{immediately} observable.


    \item \textbf{Phase}. Compare the two states:
    \begin{equation*}
        \frac{\ket{0} + \ket{1}}{\sqrt{2}}, \quad \frac{\ket{0} + i\ket{1}}{\sqrt{2}}
    \end{equation*}
    The amplitudes have identical magnitudes $\frac{1}{\sqrt{2}}$. Yet the two quantum states are different because their \textbf{relative phases} are different. The phase is part of the quantum state, but it is \textbf{not directly visible through a computational-basis measurement}.
\end{enumerate}
\textcolor{Green3}{\faIcon{question-circle} \textbf{Then why is phase useful?}} At this point it may seem that phases are useless. After all, if they cannot be measured, \emph{why do quantum algorithms spend so much effort creating them?} The answer is that \textbf{phases become observable when different branches of a superposition interfere}. We do not measure the phase itself. Instead, we first \textbf{convert the phase into amplitudes}, and \textbf{then} measure those amplitudes.

\highspace
\begin{flushleft}
    \textcolor{Green3}{\faIcon{search} \textbf{The Hadamard gate as a phase detector}}
\end{flushleft}
The most \hl{common gate used to convert phase information into amplitude information is the Hadamard gate.} Recall its action:
\begin{equation*}
    H\ket{+} = \ket{0}, \quad H\ket{-} = \ket{1}
\end{equation*}
Before the Hadamard, the states:
\begin{equation*}
    \ket{+} = \frac{\ket{0} + \ket{1}}{\sqrt{2}}, \quad \ket{-} = \frac{\ket{0} - \ket{1}}{\sqrt{2}}
\end{equation*}
Both have equal amplitudes:
\begin{equation*}
    P(0) = P(1) = \frac{1}{2}
\end{equation*}
A direct measurement cannot distinguish them. After applying the Hadamard, the states become:
\begin{equation*}
    \ket{+} \xrightarrow{H} \ket{0}, \quad \ket{-} \xrightarrow{H} \ket{1}
\end{equation*}
Now the measurement probabilities are:
\begin{equation*}
    P(0) = 1, \quad P(1) = 0
\end{equation*}
For $\ket{+}$, and:
\begin{equation*}
    P(0) = 0, \quad P(1) = 1
\end{equation*}
For $\ket{-}$. The Hadamard has converted \textbf{relative phase} into \textbf{observable probability}.

\highspace
\begin{examplebox}[: Making the Kicked-Back Phase Observable]
    Suppose phase kickback produces:
    \begin{equation*}
        \frac{\ket{0} - \ket{1}}{\sqrt{2}}
    \end{equation*}
    Suppose phase kickback produces:
    \begin{equation*}
        P(0) = P(1) = \frac{1}{2}
    \end{equation*}
    No information is revealed. Now apply a Hadamard:
    \begin{equation*}
        H\left(\frac{\ket{0} - \ket{1}}{\sqrt{2}}\right) = \ket{1}
    \end{equation*}
    A measurement now yields:
    \begin{equation*}
        P(0) = 0, \quad P(1) = 1
    \end{equation*}
    The phase has become completely observable. Not because we measured it directly, but because interference transformed it into an amplitude difference.
\end{examplebox}